\section{Calculation of the coefficient}\label{sec:numeval}
Formula~\eqref{eq:main} allows us to calculate $\cerr$ from the physical profile. For each value of
$x$ the sixteen contributions to $d(x)$ must be summed over the positions on the board. The first
evaluations, which averaged a few central squares, gave spatial errors of $10^{-5}$ to $10^{-3}$.
In particular, the line resolvent had to be evaluated again for configurations that shared a square.
The rectangle method sums all the centers of an incidence pattern on the same grid. Once the values
at $x$ have been obtained, one integrates in that variable and adds the boundary terms.

\subsection{The rectangle method}\label{sec:rectangulos}
The idea comes from the rectangle technique for the constant term~\cite[\S4]{tc}. There the norm
of an error matrix $E$ is controlled by
\[
 \operatorname{tr}\bigl((EE^{\mathsf T})^2\bigr)=\sum_{i,i',j,j'}E_{ij}E_{i'j}E_{i'j'}E_{ij'},
\]
a sum over rectangles with two rows and two columns. Each summand is expanded into sixteen
patterns $T\subseteq\{1,2,3,4\}$, according to which vertices carry their diagonal lines, with sign
$(-1)^{4-|T|}$. Since $1-4+6-4+1=0$, the leading order cancels exactly. The method carries that
bookkeeping over to the sixteen contributions to $d(x)$. Each is a sum over short closed walks of
squares joined by lines, triangles, and walks of up to four vertices with the four families, and
the patterns are grouped before integration. Thus each group is the trace of a product of small
matrices on the same grid.

The accuracy needed is dictated by the endpoint $x=0$, where $k/x^2$ amplifies the errors. With a
sixteen-point Gauss--Legendre rule on $[0,1]$, the error in one row at the first node enters the
integral multiplied by $w_1/x_1^2\approx483$. For this reason a series is used near the origin.
The numerical calculation aims for an accuracy of order $10^{-10}$ at the nodes, and the derivatives
at the origin are also evaluated. The rearrangement of the sums is exact on each grid. Separately,
with a grid profile with exact marginals, the residuals of the two cancellations were measured in
floating point. For $5\le m\le48$, in both parities, the observed values satisfied
$|k_m(0)|\le7.5\cdot10^{-13}$ and $|d_m'(0)|\le1.3\cdot10^{-11}$ \status{N}.

\subsection{A grid closed under line intersections}
We number the rows and columns starting at zero and place the square $(i,j)$ at
$((i+1/2)/m,(j+1/2)/m)$. There are $6m-2$ lines. The diagonals with integer indices $i-j$ and $i+j$
that pass through the square have the same parity. Conversely, two diagonal indices of the same
parity give a square if their intersection falls inside the board. In the continuum the two parities
give four oblique sheets, in addition to the axial sheets $R,C$. Let $\nu_1$ be the profile density
of a complete solution, denoted by $\nu$ in Appendix~\ref{app:predictor}.
For $x>0$ we take $U_f=u_f+x(1-u_f)$ on an oblique line and
$W=\nu_1/(U_DU_A)$. When lines $a,b$ intersect, $K_{ab}$ is the value of $\nu_1$ at the intersection
multiplied by $x/U_f$ for each oblique line $f$ of the pair. The other entries vanish. The diagonal
matrix $\rho$ has value $1/m$ on the axial lines and $U_f/(xm)$ on the oblique lines. Let $v$ be
the vector with values $1,-1,0$ on rows, columns, and oblique lines, respectively. The kernel of the
contractions is
\[
 \Pi=vv^{\mathsf T}/2,\qquad
 L_0=K(I+\rho K+\rho\Pi)^{-1}\rho K-\tfrac32vv^{\mathsf T}.
\]
The integrand at $x=0$ is defined by continuation. Both sheets are retained on the grid when an
observable takes values on a single parity.

For a square weight $h$, let $N_h(a,b)=h(a\cap b)$ if the distinct lines $a,b$ intersect in a grid
square. The other entries are zero. A product of these matrices sums walks along line incidences.
The trace identifies the last index with the first, and the diagonal matrices insert the line
measures. For example, $\operatorname{tr}(N_1D_1N_2D_2N_3D_3)$ sums the product of the three
weights over the triangles. Figure~\ref{fig:rectangulos} shows two rectangles in the part of
$\operatorname{tr}(N_{h_1}N_{h_2}N_{h_3}N_{h_4})$ that alternates rows and columns. The sixteen
contributions to $d$ are evaluated by matrix products, bilinear forms, and local sums. Walks with
up to four vertices appear in row $f_6$. The trace rearranges a finite sum exactly. Comparing it
with a direct count on the same grid checks the implementation. To estimate the continuum integral
one must vary $m$.
\begin{figure}[htbp]
\centering
\begin{tikzpicture}[x=.56cm,y=.56cm,
  vertice/.style={circle,draw=btinta,fill=white,text=btinta,font=\scriptsize\bfseries,
    inner sep=0pt,minimum size=3.5mm},
  paso/.style={->,>=stealth,draw=btinta,line width=.85pt}]
  \foreach \dx/\izq/\der in {0/1/3,9.7/2/4}{
    \begin{scope}[shift={(\dx,0)}]
      \bcuadriculacero{6}
      \node[vertice] (v1) at ({\izq+.5},4.5) {1};
      \node[vertice] (v2) at ({\izq+.5},1.5) {2};
      \node[vertice] (v3) at ({\der+.5},1.5) {3};
      \node[vertice] (v4) at ({\der+.5},4.5) {4};
      \draw[paso] (v1)--(v2);
      \draw[paso] (v2)--(v3);
      \draw[paso] (v3)--(v4);
      \draw[paso] (v4)--(v1);
      \node[font=\small] at (3,-.73)
        {$R_1\to C_{\izq}\to R_4\to C_{\der}\to R_1$};
    \end{scope}
  }
\end{tikzpicture}
\caption{Two summands of a trace of four incidences on the $6\times6$ grid. Vertex $k$ marks the
intersection where $N_{h_k}$ evaluates $h_k$. On the left, the product is
$h_1(1,1)h_2(4,1)h_3(4,3)h_4(1,3)$. On the right, each column increases by one.
The fourth factor returns to the initial row $R_1$. The sum runs over ordered line indices.
Each formula specifies the measures and multiplicities.}\label{fig:rectangulos}
\end{figure}


\subsection{From the grid sum to the value at each node}
At accidental coincidences the limiting value of the continuum is retained. In the sum over squares
a half weight is assigned to the side of length zero. Changing $m$ also changes the position of the
folds relative to the nodes. Thus grids of the same congruence class are taken for each fit in even
powers of $1/m$, as in Richardson extrapolation~\cite{richardson}. If there are enough new grids for
$Q_2$, the classes $0$ and $12$ modulo $24$ are fitted separately, with powers $2,4,6$. The mean of
the two intercepts is taken. Their difference, a change in the minimum grid, and the largest fit
residual measure the sensitivity. The resulting bars have status [N]. Certifying this extrapolation
would require bounding the spatial remainder in each class.

For the boundary sum the numerical analysis gives
\[
 J=-8.2386091222\ \pm5.0\,10^{-9},\qquad
 P_0=-7.7933777386\ \pm5.0\,10^{-9}\qquad\text{[N]}.
\]
For $J$ one orientation of the pair is summed and the other is included by symmetry. Adding
$P_T=0.4452313836$ gives $P_0$. The square interaction enters separately with a negative sign,
with $S=15.6334582804$. Section~\ref{sec:numcert} gives the rigorous enclosures of $P_T$ and $S$.

\subsection{The singular endpoint and the error budget}
To integrate with the weight $x^{-2}$ it is convenient first to extract the part with an explicit
antiderivative. We separate the rook functional by setting
\[
 k_F=F-F(0)-xF',\qquad
 k_6=k-k_F=d+\mathcal V-x\mathcal V'+F(0)-1/12.
\]
Differentiating $(F(x)-F(0))/x$ shows that the integral of $k_F/x^2$ equals $F(0)+F'(0)-F(1)$.
In $k_6$ the cancellations of the constant and linear terms leave a series that can be integrated
down to zero. It is reconstructed from 32 equally spaced points on $|x|=0.03$ (17 independent by
conjugation). Six Gauss--Legendre nodes on $[0.03;0.06]$ and twenty on $[0.06;1]$ are then used.
In the budget below, the errors in the nodal data are estimated by varying the spatial grid. The
series and quadrature bound follows from the assumed tolerances for the contour samples in
Appendix~\ref{app:numquad}.

\begin{center}\small
\begin{tabular}{lr}
Source & Contribution to the numerical error budget\\ \hline
Nodal data from the twenty-point panel & $2.8\,10^{-8}$\\
Data from the circle of radius $0.03$ & $1.4\,10^{-8}$\\
Boundary value $P_0$ & $5.0\,10^{-9}$\\
Nodal data from the six-point panel & $2.3\,10^{-9}$\\
Series and quadrature, conditional upper bound & $8.74\,10^{-10}$\\
Value of $S$ in the numerical calculation & $2.0\,10^{-12}$
\end{tabular}
\end{center}
The calculation gives $0.4055003786$. The largest contributions to the error come from the values
on the twenty-point panel and on the circle. The quadratic combination of the indicators is
$3.2\,10^{-8}$ and the linear one is $5.0\,10^{-8}$. When the fit for $Q_2$ was changed, the result
varied between $-3.7\,10^{-9}$ and $+7.6\,10^{-9}$ relative to the reference. On the basis of these
empirical checks we report $0.40550038\pm3\,10^{-8}$ [N], with the conservative estimate
$0.4055004\pm2\,10^{-7}$ [N].
